% ============================================================================
%  Problems: ลำดับ (Sequences)
%  Ordered from easy (basic) → hard (POSN level).
%  Shared by exercise.tex and answer-key.tex
% ============================================================================

% --- Part 1: พื้นฐาน (Basic) ---
\examsection{ตอนที่ 1: พื้นฐาน — รู้จักลำดับ}

\problem{
  จงระบุว่าแต่ละลำดับต่อไปนี้เป็นลำดับเลขคณิตหรือลำดับเรขาคณิต
  พร้อมทั้งหา $d$ หรือ $r$
  \begin{enumerate}
    \item[(ก)] $5, 10, 15, 20, 25, \ldots$
    \item[(ข)] $2, 4, 8, 16, 32, \ldots$
    \item[(ค)] $100, 90, 80, 70, \ldots$
  \end{enumerate}
}{
  (ก) เลขคณิต, $d = 5$

  (ข) เรขาคณิต, $r = 2$

  (ค) เลขคณิต, $d = -10$
}

\problem{
  จงเขียน 4 พจน์แรกของลำดับที่กำหนดให้
  \begin{enumerate}
    \item[(ก)] $a_n = 2n + 1$
    \item[(ข)] $a_n = 3^n$
    \item[(ค)] $a_n = n^2 - n$
  \end{enumerate}
}{
  (ก) $3, 5, 7, 9$

  (ข) $3, 9, 27, 81$

  (ค) $0, 2, 6, 12$
}

\problem{
  จงหาค่าของ $d$ (ผลต่างร่วม) สำหรับลำดับเลขคณิตต่อไปนี้
  \begin{enumerate}
    \item[(ก)] $7, 12, 17, 22, \ldots$
    \item[(ข)] $20, 14, 8, 2, \ldots$
  \end{enumerate}
}{
  (ก) $d = 5$

  (ข) $d = -6$
}

\problem{
  จงหาค่าของ $r$ (อัตราส่วนร่วม) สำหรับลำดับเรขาคณิตต่อไปนี้
  \begin{enumerate}
    \item[(ก)] $5, 15, 45, 135, \ldots$
    \item[(ข)] $64, 32, 16, 8, \ldots$
  \end{enumerate}
}{
  (ก) $r = 3$

  (ข) $r = \frac{1}{2}$
}

\problem{
  จงหาพจน์ถัดไปอีก 2 พจน์ของลำดับฟิโบนักชี
  $1, 1, 2, 3, 5, 8, 13, \ldots$

  (Hint: แต่ละพจน์เกิดจากผลบวกของสองพจน์ก่อนหน้า)
}{
  $21, 34$
}

% --- Part 2: การประยุกต์ (Intermediate) ---
\examsection{ตอนที่ 2: การประยุกต์ — พจน์ทั่วไปและผลบวก}

\problem{
  จงหาพจน์ที่ 15 ของลำดับเลขคณิตที่มี $a_1 = 3$ และ $d = 4$
}{
  $a_{15} = 3 + (15-1) \times 4 = 3 + 56 = 59$
}

\problem{
  จงหาพจน์ที่ 30 ของลำดับเลขคณิต $5, 9, 13, 17, \ldots$
}{
  $a_1 = 5$, $d = 4$

  $a_{30} = 5 + (30-1) \times 4 = 5 + 116 = 121$
}

\problem{
  จงหาผลบวก 20 พจน์แรกของลำดับเลขคณิต $2, 5, 8, 11, \ldots$
}{
  $a_1 = 2$, $d = 3$

  $a_{20} = 2 + 19 \times 3 = 59$

  $S_{20} = \frac{20}{2}(2 + 59) = 10 \times 61 = 610$
}

\problem{
  จงหาพจน์ที่ 8 ของลำดับเรขาคณิตที่มี $a_1 = 2$ และ $r = 3$
}{
  $a_8 = 2 \times 3^7 = 2 \times 2187 = 4374$
}

\problem{
  จงหาพจน์ที่ 12 ของลำดับเลขคณิตที่มี $a_3 = 7$ และ $a_7 = 19$
  (Hint: หา $d$ จาก $a_3$ และ $a_7$ ก่อน)
}{
  $a_7 = a_3 + 4d$ \quad ดังนั้น $19 = 7 + 4d$ \quad $\Rightarrow d = 3$

  $a_3 = a_1 + 2d$ \quad ดังนั้น $7 = a_1 + 6$ \quad $\Rightarrow a_1 = 1$

  $a_{12} = 1 + 11 \times 3 = 34$
}

% --- Part 3: ระดับ สอวน. (POSN Level) ---
\examsection{ตอนที่ 3: ระดับข้อสอบ สอวน.}

\problem{
  จงหาพจน์ที่ 7 ของลำดับฟิโบนักชี $F_n$
}{
  $1, 1, 2, 3, 5, 8, 13$

  $F_7 = 13$
}

\problem{
  ลำดับเรขาคณิตหนึ่งมี $a_3 = 12$ และ $a_6 = 96$
  จงหา $a_1$ และ $r$
}{
  $a_6 = a_3 \times r^3$ \quad ดังนั้น $96 = 12 \times r^3$ \quad $\Rightarrow r^3 = 8$ \quad $\Rightarrow r = 2$

  $a_3 = a_1 \times r^2$ \quad ดังนั้น $12 = a_1 \times 4$ \quad $\Rightarrow a_1 = 3$
}

\problem{
  พจน์ที่ 8 และพจน์ที่ 10 ของลำดับเลขคณิต
  เป็น 36 และ 46 ตามลำดับ จงหาพจน์ที่ 17
}{
  $d = \frac{a_{10} - a_8}{2} = \frac{46 - 36}{2} = 5$

  $a_8 = a_1 + 7d$ \quad ดังนั้น $36 = a_1 + 35$ \quad $\Rightarrow a_1 = 1$

  $a_{17} = 1 + 16 \times 5 = 81$
}

\problem{
  ถ้า $F_n$ แทนพจน์ที่ $n$ ของลำดับฟิโบนักชี
  โดยที่ $F_1 = 1$, $F_2 = 1$
  จงหาค่าของ $F_1 + F_2 + F_3 + \cdots + F_6$
}{
  $1 + 1 + 2 + 3 + 5 + 8 = 20$
}
